\paragraph{原始综合征、单素数模审计与离线最短公钥量纲：Hankel 低秩的可核验接口}\label{par:xi-hankel-syndrome-offline-audit-optimality}
本段把 Hankel 低秩所释放的核向量视为\emph{综合征}对象，并给出三类互补的可核验接口：
其一在整数域上由格点 Siegel 引理给出原始综合征的显式范数界；
其二给出单素数模上的零误差审计机制与位复杂度上界；
其三在有限域成员判定语义下给出验证参数的不可约信息论下界与近最短实现。
此外，线性递推约束与生成函数运算的闭合律（Cauchy 卷积、Hadamard 乘积）在综合征多项式层面可由结果式完全刻画。

\begin{theorem}[Hankel 低秩的原始综合征与显式 \(\ell_\infty\) 界]\label{thm:xi-hankel-primitive-syndrome-siegel-bound}
\leanverified{paper\_xi\_hankel\_primitive\_syndrome\_siegel\_bound}
设 \(h\in\NN\)，令 \(n:=h+1\)。
给定整数序列 \((a_k)_{k=0}^{2h}\subset\ZZ\)，定义 Hankel 矩阵
\[
H:=H_h(a)=(a_{i+j})_{0\le i,j\le h}\in\ZZ^{n\times n}.
\]
若 \(\mathrm{rank}_{\QQ}(H)=\rho<n\)，令 \(\delta:=n-\rho\ge 1\)，并假设 \(|a_k|\le B\)（\(0\le k\le 2h\)）。
则存在非零整数向量 \(s\in\ker_{\QQ}(H)\cap\ZZ^n\) 使得 \(\gcd(s_0,\dots,s_{n-1})=1\) 且
\[
|s|_\infty\ \le\ (nB)^{\rho/\delta}.
\]
当 \(\delta=1\) 时，\(\ker_{\QQ}(H)\cap\ZZ^n\) 作为秩 \(1\) 的 \(\ZZ\)-模由唯一的原始生成元 \(\pm s^\ast\) 生成，并可取
\[
|s^\ast|_\infty\ \le\ n^{(n-1)/2}\,B^{\,n-1}.
\]
\end{theorem}
\begin{proof}
取 \(H\) 的 \(\rho\) 行构成的满秩子矩阵 \(A\in\ZZ^{\rho\times n}\)，则 \(\ker_{\QQ}(A)=\ker_{\QQ}(H)\)。
由格点版 Siegel 引理（或等价的 Minkowski 体积论证），存在非零 \(x\in\ZZ^n\) 使 \(Ax=0\) 且
\(
|x|_\infty\le (nB)^{\rho/(n-\rho)}=(nB)^{\rho/\delta}
\)。
令 \(s:=x/\gcd(x_0,\dots,x_{n-1})\) 得原始向量，且不增大 \(|\cdot|_\infty\)。
\par
当 \(\delta=1\) 时，\(\mathrm{rank}_{\QQ}(H)=n-1\) 且 \(\det(H)=0\)。
由恒等式 \(H\operatorname{adj}(H)=\det(H)I=0\)，\(\operatorname{adj}(H)\) 的任一非零列向量属于 \(\ker(H)\cap\ZZ^n\)。
其坐标为 \((n-1)\times(n-1)\) 余子式行列式；由 Hadamard 不等式，任一余子式的绝对值不超过
\(
n^{(n-1)/2}B^{n-1}
\)。
核维数为 \(1\) 保证原始生成元唯一（至多差符号）。证毕。
\end{proof}

\begin{theorem}[余子式综合征与单坐标判别]\label{thm:xi-hankel-cofactor-syndrome-single-coordinate}
\leanverified{paper\_xi\_hankel\_cofactor\_syndrome\_single\_coordinate}
设 \(R\) 为交换整环，给定整数 \(d\ge 1\) 与序列 \((a_k)_{k=0}^{2d}\subset R\)。
令
\[
H:=H_d(a):=(a_{i+j})_{0\le i,j\le d}\in R^{(d+1)\times(d+1)},
\qquad
H_0:=(a_{i+j})_{0\le i,j\le d-1}\in R^{d\times d}.
\]
假设在分式域 \(K:=\mathrm{Frac}(R)\) 上有 \(\rank_K(H)=d\)（等价地：\(\det(H)=0\) 且 \(\ker_K(H)\) 为一维）。
定义向量
\[
u=(u_0,\dots,u_d)^{\mathsf T}\in R^{d+1},
\qquad
u_j:=(-1)^{d+j}\det\bigl(H^{(d,j)}\bigr),
\]
其中 \(H^{(d,j)}\) 表示从 \(H\) 删去第 \(d\) 行与第 \(j\) 列所得的 \(d\times d\) 子矩阵。
则：
\begin{enumerate}
  \item \(u\neq 0\) 且 \(Hu=0\)，从而 \(u\) 给出 Hankel 低秩关系的一组可核验综合征；
  \item 末坐标满足 \(u_d=\det(H_0)\)；
  \item 令 \(s^\ast:=u/\gcd(u_0,\dots,u_d)\in R^{d+1}\) 为原始综合征（按整除意义），则 \(s^\ast\) 在 \(\ker(H)\cap R^{d+1}\) 中按单位倍唯一，并且 \(s^\ast_d\mid \det(H_0)\)。
\end{enumerate}
\end{theorem}
\begin{proof}
由伴随矩阵恒等式
\(
H\operatorname{adj}(H)=\operatorname{adj}(H)H=\det(H)\,I_{d+1}
 \)
并用 \(\det(H)=0\) 可知 \(\operatorname{adj}(H)\) 的任一列属于 \(\ker(H)\subset K^{d+1}\)。
在 \(\rank_K(H)=d\) 下，存在某个 \(d\times d\) 余子式在 \(K\) 中非零，故 \(\operatorname{adj}(H)\neq 0\)。
取其第 \(d\) 列得到非零核向量；按 \(\operatorname{adj}(H)\) 的定义，该列坐标恰为上式所给 \(u\)，从而 \(Hu=0\) 且 \(u\neq 0\)。
又 \(u_d\) 对应删去第 \(d\) 行与第 \(d\) 列的余子式，故 \(u_d=\det(H_0)\)。
原始化与整除结论由定义立即推出；并可与 Toeplitz--扫描 Hankel 的合同分解口径相对照（\cite{Internal2026ToeplitzScanHankelBlaschkeVanVleck}）。证毕。
\end{proof}

\begin{corollary}[最大子式综合征的高度界与确定性 CRT 复原]\label{cor:xi-hankel-maxminor-height-bound-deterministic-crt-reconstruction}
在定理 \ref{thm:xi-hankel-cofactor-syndrome-single-coordinate} 的设定下取 \(R=\ZZ\)，并进一步设
\[
H_0\in\ZZ^{d\times d}\ \ \text{可逆于}\ \QQ,\qquad
\widehat H\in\ZZ^{d\times(d+1)}\ \ \text{如}\ \ref{thm:xi-hankel-maximal-minor-syndrome-normal-form-uniqueness}\ \text{所定义}.
\]
令 \(\Delta=(\Delta_0,\dots,\Delta_d)^{\mathsf T}\in\ZZ^{d+1}\) 为由 \(\widehat H\) 的最大子式给出的综合征向量（定理 \ref{thm:xi-hankel-maximal-minor-syndrome-normal-form-uniqueness}），并记
\[
A:=\max_{0\le n\le 2d-1}|a_{n+\ell}|.
\]
则：
\begin{enumerate}
  \item \textbf{（高度界）}对一切 \(0\le k\le d\) 有
  \[
  |\Delta_k|\ \le\ d^{d/2}\,A^{d}.
  \]
  \item \textbf{（确定性 CRT 复原）}令 \(B:=d^{d/2}A^d\)。
  取两两不同素数 \(p_1,\dots,p_s\) 满足 \(p_i\nmid \det(H_0)\) 且
  \[
  P:=\prod_{i=1}^{s}p_i\ >\ 2B.
  \]
  则由每个模 \(p_i\) 下的余子式公式计算得到的 \(\Delta\bmod p_i\)，可经中国剩余定理唯一复原出满足 \(|\Delta_k|\le B\) 的整数向量 \(\Delta\)（并自动满足 \(\Delta_d=\det(H_0)\)）。
\end{enumerate}
\end{corollary}
\begin{proof}
(1) 由定理 \ref{thm:xi-hankel-maximal-minor-syndrome-normal-form-uniqueness}，每个 \(\Delta_k\) 是某个 \(d\times d\) 子式行列式。
该子式的每行欧氏范数不超过 \(\sqrt d\,A\)，由 Hadamard 不等式得
\(
|\Delta_k|\le (\sqrt d\,A)^d=d^{d/2}A^d
\)。
\par
(2) 条件 \(p_i\nmid \det(H_0)\) 确保模 \(p_i\) 下 \(H_0\) 可逆，从而最大子式综合征在 \(\FF_{p_i}\) 上不退化并由同一余子式公式无歧义计算。
由同余数据与 CRT 得到唯一的 \(\Delta\bmod P\)。
由于 \(P>2B\)，每个坐标 \(\Delta_k\) 在同余类中存在唯一代表落在区间 \((-\tfrac{P}{2},\tfrac{P}{2})\)，并且该代表必满足 \(|\Delta_k|\le B\)，从而得到整数层面的唯一复原。
最后，\(\Delta_d=\det(H_0)\) 在整数上成立（定理 \ref{thm:xi-hankel-maximal-minor-syndrome-normal-form-uniqueness}），故复原向量亦满足该一致性。证毕。
\end{proof}

\begin{corollary}[坏素数的单坐标判别与总刚性读出]\label{cor:xi-hankel-badprimes-single-coordinate-stiffness}
在定理 \ref{thm:xi-hankel-cofactor-syndrome-single-coordinate} 的设定下，进一步取 \(R=\ZZ\)。
对任意素数 \(p\)，下列条件等价：
\[
H_0\bmod p\ \text{可逆}
\quad\Longleftrightarrow\quad
p\nmid u_d
\quad\Longleftrightarrow\quad
p\nmid \det(H_0).
\]
因此当 \(\ker_{\QQ}(H)\) 为一维时，导致模 \(p\) 下 Hankel 主块退化的坏素数集合完全由单个坐标 \(u_d\) 的素因子集合决定。
并且对任意素数 \(p\) 有
\[
v_p(\det(H_0))=v_p(u_d),
\]
其与命题 \ref{prop:xi-hankel-det-reach-observe-factorization-padic} 的 \(p\)-进总刚性 \(v_p(\det H_0)\) 一致。
\end{corollary}
\begin{proof}
由定理 \ref{thm:xi-hankel-cofactor-syndrome-single-coordinate}，有 \(u_d=\det(H_0)\)，故上述等价性与赋值恒等式均为直接代入。证毕。
\end{proof}

\begin{proposition}[单素数模审计的零误差化与位复杂度上界]\label{prop:xi-hankel-single-prime-audit-zero-error}
\leanverified{paper\_xi\_hankel\_single\_prime\_audit\_zero\_error}
沿用定理 \ref{thm:xi-hankel-primitive-syndrome-siegel-bound} 的记号。
给定任意整数向量 \(s\in\ZZ^n\)，令
\[
v:=Hs\in\ZZ^n,\qquad M:=n^2B\,|s|_\infty.
\]
则对任意素数 \(p>M\)，有等价性
\[
Hs\equiv 0\pmod p\quad\Longleftrightarrow\quad Hs=0\ \text{于}\ \ZZ^n.
\]
并且当 \(M\ge 2\) 时存在素数 \(p\) 满足 \(M<p<2M\)，从而可取
\[
\log_2 p\ \le\ \log_2(2M)=O\!\bigl(\log n+\log B+\log|s|_\infty\bigr).
\]
\end{proposition}
\begin{proof}
对每个分量 \(v_i=\sum_{j=0}^{n-1}a_{i+j}s_j\)，有
\[
|v_i|
\le \sum_{j=0}^{n-1}|a_{i+j}|\,|s_j|
\le nB\,|s|_1
\le n^2B\,|s|_\infty
=M.
\]
若 \(p>M\) 且 \(v\neq 0\)，则存在 \(i\) 使 \(0<|v_i|<p\)，从而 \(v_i\not\equiv 0\pmod p\)。
反向方向显然。
素数存在性由 Bertrand--Chebyshev 定理给出。证毕。
\end{proof}

\begin{theorem}[多素数离线审计：以乘积模数替代单一大素数]\label{thm:xi-hankel-multiprime-offline-audit}
\leanverified{paper\_xi\_hankel\_multiprime\_offline\_audit}
沿用命题 \ref{prop:xi-hankel-single-prime-audit-zero-error} 的记号，并令
\(
v:=Hs\in\ZZ^n
\)、
\(
M:=n^2B\,|s|_\infty
\)。
取两两不同素数 \(q_1,\dots,q_K\) 并记
\[
Q:=\prod_{i=1}^{K}q_i.
\]
若 \(Q>2M\)，则下列命题等价：
\begin{enumerate}
  \item \(Hs=0\)（在整数上）；
  \item \(Hs\equiv 0\pmod{q_i}\) 对所有 \(i=1,\dots,K\) 成立。
\end{enumerate}
\end{theorem}
\begin{proof}
\((1)\Rightarrow(2)\) 平凡。
若 \((2)\) 成立，则 \(v\equiv 0\pmod{q_i}\) 对所有 \(i\) 成立，从而 \(v\equiv 0\pmod Q\)。
由命题 \ref{prop:xi-hankel-single-prime-audit-zero-error} 的估计，对每个分量 \(v_j\) 有 \(|v_j|\le M<Q/2\)。
而整数满足 \(v_j\equiv 0\pmod Q\) 且 \(|v_j|<Q/2\) 时只能 \(v_j=0\)。
故 \(v=0\)，即 \(Hs=0\)。证毕。
\end{proof}

\begin{corollary}[最大模素数受限时的最优对数尺度]\label{cor:xi-hankel-offline-audit-max-prime-logscale}
在定理 \ref{thm:xi-hankel-multiprime-offline-audit} 的设定下，若以“最小化最大模素数 \(\max_i q_i\)”为设计准则并保持素数互异，则最优值 \(q_{\max}^\ast\) 由 Chebyshev 函数
\(
\vartheta(x)=\sum_{p\le x}\log p
\)
刻画为
\[
q_{\max}^\ast=\inf\{x\ge 2:\ \vartheta(x)\ge \log(2M)\}.
\]
因此当 \(M\to\infty\) 时有渐近
\[
q_{\max}^\ast\sim \log(2M),
\qquad
K^\ast\sim \frac{\log(2M)}{\log\log(2M)},
\]
其中 \(K^\ast\) 为达到 \(Q>2M\) 所需的最小素数轴数量。
\end{corollary}
\begin{proof}
由于要满足 \(Q>2M\)，必要条件是 \(\sum_i\log q_i>\log(2M)\)。
在素数互异且最小化最大值的准则下，取不超过某个 \(x\) 的全部素数能最大化 \(\sum\log p\)，其上界正是 \(\vartheta(x)\)，故最优最大素数满足所示刻画。
当 \(M\to\infty\) 时，由 \(\vartheta(x)\sim x\) 得 \(q_{\max}^\ast\sim \log(2M)\)。
又 \(K^\ast=\pi(q_{\max}^\ast)\)，由 \(\pi(x)\sim x/\log x\) 得 \(K^\ast\sim \log(2M)/\log\log(2M)\)。证毕。
\end{proof}

% Multi-stage prime-axis allocation under disjointness constraint.

\begin{theorem}[多阶段 CRT 素数轴分配的加性控制律]\label{thm:xi-crt-multistage-axis-allocation-additive-control}
设 \(k\ge 1\)，并给定阈值 \(A_1,\dots,A_k>1\)。
要求选择两两不交的素数集合 \(S_1,\dots,S_k\subset\mathbb P\) 使得
\[
\prod_{p\in S_i}p\ \ge\ A_i\qquad(1\le i\le k),
\]
并最小化全局最大素数
\[
p_{\max}:=\max\Bigl(\bigcup_{i=1}^k S_i\Bigr).
\]
记 Chebyshev 函数
\[
\vartheta(x):=\sum_{p\le x}\log p.
\]
则最优值满足精确刻画
\[
\boxed{\;
p_{\max}^\ast
=\inf\left\{x\ge 2:\ \vartheta(x)\ge \sum_{i=1}^k\log A_i\right\}. \;}
\]
并且当 \(\sum_i\log A_i\to\infty\) 时有渐近
\[
p_{\max}^\ast\ \sim\ \sum_{i=1}^k\log A_i,
\]
其中 \(\sim\) 取在实数意义下。
\leanverified{paper\_xi\_crt\_multistage\_axis\_allocation\_additive\_control}
\end{theorem}
\begin{proof}
\textbf{必要性。}
任取可行的 \((S_i)\)。
取对数并求和得
\[
\sum_{i=1}^k\log A_i
\le
\sum_{i=1}^k\sum_{p\in S_i}\log p
=
\sum_{p\in\cup_i S_i}\log p
\le
\sum_{p\le p_{\max}}\log p
=\vartheta(p_{\max}),
\]
从而若 \(\vartheta(x)<\sum_i\log A_i\) 则任何满足 \(p_{\max}\le x\) 的方案均不可能可行，得到下确界形式的下界。
\par
\textbf{充分性。}
若 \(\vartheta(x)\ge \sum_i\log A_i\)，将全部不超过 \(x\) 的素数按从小到大排列并依次分配：
先把素数逐个加入 \(S_1\) 直至 \(\sum_{p\in S_1}\log p\ge \log A_1\)，再继续向 \(S_2\) 分配直至满足 \(\log A_2\)，依此类推。
由于总和 \(\sum_{p\le x}\log p\) 足以覆盖 \(\sum_i\log A_i\)，该过程必在用尽全部 \(\le x\) 的素数前完成 \(k\) 个阈值，因此得到一组可行解并满足 \(p_{\max}\le x\)。
故最小可能的 \(p_{\max}\) 恰为所示下确界。
\par
渐近式由素数定理等价形式 \(\vartheta(x)\sim x\) 直接推出。
证毕。
\end{proof}

\begin{corollary}[多阶段 prime-axis 数量的统一标度]\label{cor:xi-crt-multistage-axis-count-scaling}
\leanverified{paper\_xi\_crt\_multistage\_axis\_count\_scaling}
在定理 \ref{thm:xi-crt-multistage-axis-allocation-additive-control} 的设定下，令
\(
L:=\sum_{i=1}^k\log A_i
\)。
当 \(L\to\infty\) 时，达到最优最大素数标度 \(p_{\max}^\ast\sim L\) 所需的最小总素数轴数 \(K^\ast\) 满足
\[
K^\ast\ \sim\ \frac{L}{\log L}.
\]
\end{corollary}
\begin{proof}
由定理 \ref{thm:xi-crt-multistage-axis-allocation-additive-control} 得 \(p_{\max}^\ast\sim L\)。
取最小轴数时，可取全部不超过 \(p_{\max}^\ast\) 的素数（不重复），于是 \(K^\ast=\pi(p_{\max}^\ast)\)。
由素数定理 \(\pi(x)\sim x/\log x\) 得
\(
K^\ast\sim p_{\max}^\ast/\log p_{\max}^\ast\sim L/\log L
\)。
证毕。
\end{proof}

\endinput


% One-round randomized prime audit bound.

\begin{theorem}[随机素数单轮审计的误接受上界]\label{thm:xi-random-prime-single-round-audit-false-accept}
\leanverified{paper\_xi\_random\_prime\_single\_round\_audit\_false\_accept}
在“综合征约束—离线审计”语义下，令 \(\{e_{\alpha}\}_{\alpha\in\mathcal I}\subset\ZZ\) 为有限族待核验的整数约束量，并假设至少存在某个 \(\alpha\) 使 \(e_\alpha\neq 0\)。
取显式有效上界 \(U\ge \max_{\alpha\in\mathcal I}|e_\alpha|\)，并令
\[
g:=\gcd_{\alpha\in\mathcal I}(e_\alpha)\ \in\ \ZZ\setminus\{0\}.
\]
从区间 \([P,2P]\) 内的素数集合中等概率抽取一素数 \(p\)。
则误接受概率满足
\[
\Pr\bigl(\forall \alpha\in\mathcal I:\ e_\alpha\equiv 0\ (\mathrm{mod}\ p)\bigr)
\ \le\
\frac{\omega(|g|)}{\pi(2P)-\pi(P)},
\qquad
\omega(|g|)\le \frac{\log|g|}{\log 2}\le \frac{N\log U}{\log 2},
\]
其中 \(N:=|\mathcal I|\)，\(\omega(m)\) 为 \(m\) 的不同素因子个数，\(\pi(\cdot)\) 为素数计数函数。
特别地，若 \(P>U\)，则误接受概率为 \(0\)。
\end{theorem}
\begin{proof}
若在模 \(p\) 下所有约束均通过，则 \(p\mid e_\alpha\) 对一切 \(\alpha\) 成立，故 \(p\mid g\)。
因此可能导致误接受的素数 \(p\in[P,2P]\) 必属于 \(g\) 的素因子集合，其个数至多 \(\omega(|g|)\)。
等概率抽取给出第一不等式。
\par
由于每个不同素因子至少贡献一个 \(2\)，有 \(2^{\omega(|g|)}\le |g|\)，从而 \(\omega(|g|)\le \log|g|/\log 2\)。
又由 \(|g|\le \prod_{\alpha}|e_\alpha|\le U^{N}\) 得 \(\log|g|\le N\log U\)，从而得到末尾界。
若 \(P>U\)，则区间内任意素数 \(p\) 不可能整除某个满足 \(0<|e_\alpha|\le U\) 的非零约束量，因此不可能误接受。证毕。
\end{proof}

\endinput


\begin{proposition}[一次模素数秩审计的缺口恢复与失败概率]\label{prop:xi-hankel-defect-single-prime-rank-audit}
令 \(H\in\ZZ^{n\times n}\) 为整数 Hankel 矩阵，并记
\[
d:=\mathrm{rank}_{\QQ}(H),
\qquad
\delta:=n-d.
\]
则存在有限素数集 \(\mathcal B(H)\) 使得对一切 \(p\notin\mathcal B(H)\) 有
\[
\mathrm{rank}_{\FF_p}(H\bmod p)=d,
\qquad
\dim_{\FF_p}\ker(H\bmod p)=\delta.
\]
并且若从区间 \([P,2P]\) 内的素数中按均匀分布抽取 \(p\)，则
\[
\PP[p\in\mathcal B(H)]
\ \le\
\frac{|\mathcal B(H)|}{\pi(2P)-\pi(P)}
\ \xrightarrow[P\to\infty]{}\ 0.
\]
\end{proposition}

\begin{proof}
由 \(d=\mathrm{rank}_{\QQ}(H)\)，存在某个 \(d\times d\) 子式 \(\Delta\in\ZZ\setminus\{0\}\)，并且所有 \((d+1)\times(d+1)\) 子式均为零整数。
令
\(
\mathcal B(H):=\{p:\ p\mid \Delta\}
\)，则 \(\mathcal B(H)\) 有限。
若 \(p\notin\mathcal B(H)\)，则 \(\Delta\bmod p\neq 0\)，从而 \(\mathrm{rank}_{\FF_p}(H\bmod p)\ge d\)；
另一方面，由所有 \((d+1)\times(d+1)\) 子式恒为零，得 \(\mathrm{rank}_{\FF_p}(H\bmod p)\le d\)，故秩恰为 \(d\)。
于是核维数为 \(n-d=\delta\)。
\par
概率界由计数恒等式给出：区间 \([P,2P]\) 内能整除 \(\Delta\) 的素数至多为 \(|\mathcal B(H)|\) 个，除以区间内素数总数 \(\pi(2P)-\pi(P)\) 即得；
极限由 \(\pi(2P)-\pi(P)\to\infty\) 推出。证毕。
\end{proof}

\leanverified{paper\_xi\_linear\_recurrence\_local\_test\_by\_violation\_rate}
\begin{proposition}[线性递推约束的局部核验与拒绝率下界（按违反率）]\label{prop:xi-linear-recurrence-local-test-by-violation-rate}
令 \(\FF\) 为任意域，取 \(r\ge 1\) 与 \(c=(c_0,\dots,c_r)\in\FF^{r+1}\) 满足 \(c_0c_r\neq 0\)。
对任意长度 \(N>r\) 定义线性约束族
\[
\mathcal C_N(c):=\Bigl\{a=(a_0,\dots,a_{N-1})\in\FF^N:\ \sum_{j=0}^r c_j a_{i+j}=0,\ \forall\,0\le i\le N-r-1\Bigr\}.
\]
对任意 \(a\in\FF^N\)，定义违反约束的索引集合
\[
V(a):=\Bigl\{0\le i\le N-r-1:\ \sum_{j=0}^r c_j a_{i+j}\ne 0\Bigr\},
\qquad
\delta_{\mathrm{viol}}(a):=\frac{|V(a)|}{N-r}\in[0,1].
\]
若以均匀随机方式抽取 \(t\) 个索引 \(i\in\{0,\dots,N-r-1\}\) 检查相应约束，则拒绝概率至少为
\[
1-(1-\delta_{\mathrm{viol}}(a))^t.
\]
特别地，当 \(\delta_{\mathrm{viol}}(a)\ge\varepsilon>0\) 时，取
\(
t\ge \lceil \varepsilon^{-1}\log 3\rceil
\)
即可保证拒绝概率 \(\ge 2/3\)，且该抽样复杂度与 \(N\) 无关。
\end{proposition}
\begin{proof}
对单次均匀抽样，命中违反索引的概率恰为 \(\delta_{\mathrm{viol}}(a)\)。
独立重复抽样给出通过概率至多 \((1-\delta_{\mathrm{viol}}(a))^t\)，从而得到拒绝概率下界。
最后一条由 \((1-\varepsilon)^t\le 1/3\) 等价改写得到。证毕。
\end{proof}

\begin{theorem}[Cauchy 卷积在既约分母上的闭合律]\label{thm:xi-cauchy-convolution-denominator-gcd-law}
\leanverified{paper\_xi\_cauchy\_convolution\_denominator\_gcd\_law}
设 \(K\) 为特征 \(0\) 的域。
令两列序列 \((a_n)_{n\ge 0},(b_n)_{n\ge 0}\subset K\) 的生成函数分别为既约有理函数
\[
A(z)=\sum_{n\ge 0}a_n z^n=\frac{U_a(z)}{V_a(z)},\qquad
B(z)=\sum_{n\ge 0}b_n z^n=\frac{U_b(z)}{V_b(z)},
\]
其中 \(U_a,V_a,U_b,V_b\in K[z]\)、\(\gcd(U_a,V_a)=\gcd(U_b,V_b)=1\)，并假设 \(V_a(0)=V_b(0)=1\)。
定义 Cauchy 卷积 \(c_n:=\sum_{k=0}^n a_k b_{n-k}\) 及 \(C(z)=\sum_{n\ge 0}c_n z^n=A(z)B(z)\)。
则 \(C(z)\) 的既约分母满足精确公式
\[
V_c(z)\ =\ \frac{V_a(z)V_b(z)}{\gcd\!\bigl(V_a(z),U_b(z)\bigr)\cdot \gcd\!\bigl(V_b(z),U_a(z)\bigr)}.
\]
特别地，若 \(\gcd(V_a,U_b)=\gcd(V_b,U_a)=1\)，则 \(V_c=V_aV_b\)，从而卷积后的递推阶数严格加和（等于 \(\deg V_a+\deg V_b\)）。
\end{theorem}
\begin{proof}
由 \(C(z)=U_a(z)U_b(z)/(V_a(z)V_b(z))\)。
约分仅可能来自分母因子与分子共享的不可约因子。
由于 \(\gcd(U_a,V_a)=1\)，\(V_a\) 的任一不可约因子不可能与 \(U_a\) 共有；同理 \(V_b\) 不与 \(U_b\) 共有。
因此一切可约分因子只能来自交叉项 \(\gcd(V_a,U_b)\) 与 \(\gcd(V_b,U_a)\)，并分别整除 \(V_a\)、\(V_b\)，从而互不干扰。
故既约分母为所示商。证毕。
\end{proof}

\begin{theorem}[Hadamard 乘积的综合征多项式与结果式闭式]\label{thm:xi-hadamard-product-resultant-multiplicative-convolution}
\leanverified{paper\_xi\_hadamard\_product\_resultant\_multiplicative\_convolution}
设 \(K\) 为代数闭域。
令两列序列 \((a_n)_{n\ge 0},(b_n)_{n\ge 0}\subset K\) 具有有限指数表示
\[
a_n=\sum_{i=1}^{r_a}\alpha_i\lambda_i^{n},\qquad
b_n=\sum_{j=1}^{r_b}\beta_j\mu_j^{n},
\]
其中 \(\alpha_i,\beta_j\neq 0\)，且 \(\lambda_i,\mu_j\in K^\times\) 两两互异。
定义 Hadamard 乘积序列 \(h_n:=a_n b_n\)。
则 \(h\) 仍为有限指数表示，并满足
\[
h_n=\sum_{i=1}^{r_a}\sum_{j=1}^{r_b}(\alpha_i\beta_j)\,(\lambda_i\mu_j)^{n}.
\]
令
\[
P_a(t):=\prod_{i=1}^{r_a}(t-\lambda_i),\qquad
P_b(t):=\prod_{j=1}^{r_b}(t-\mu_j),
\]
并定义
\[
P_h(t)\ :=\ \operatorname{Res}_{x}\!\Bigl(P_a(x),\ x^{r_b}P_b\!\bigl(t/x\bigr)\Bigr)\in K[t].
\]
则恒等式
\[
P_h(t)=\prod_{i=1}^{r_a}\prod_{j=1}^{r_b}(t-\lambda_i\mu_j)
\]
成立；从而 \(P_h\) 消去序列 \(h\)。
若进一步假设所有乘积 \(\lambda_i\mu_j\) 两两互异，则 \(P_h\) 为 \(h\) 的最小特征多项式，递推阶数达到 \(r_ar_b\)。
\end{theorem}
\begin{proof}
由指数展开立得
\(
h_n=\sum_{i,j}(\alpha_i\beta_j)(\lambda_i\mu_j)^n
\)，从而任一以 \(\{\lambda_i\mu_j\}\) 为根集的多项式必消去 \(h\)。
\par
对结果式，记
\[
Q_t(x):=x^{r_b}P_b(t/x)=\prod_{j=1}^{r_b}(t-\mu_j x)\in K[t][x].
\]
由于 \(P_a\) 首一，结果式的根值表述给出
\[
\operatorname{Res}_x(P_a(x),Q_t(x))=\prod_{i=1}^{r_a}Q_t(\lambda_i)
=\prod_{i=1}^{r_a}\prod_{j=1}^{r_b}(t-\mu_j\lambda_i),
\]
即所示闭式。
当乘积根互异时，指数基底的 Vandermonde 非退化性给出最小性。证毕。
\end{proof}

\begin{theorem}[离线成员判定的不可约信息论下界（Grassmann 计数）]\label{thm:xi-offline-membership-grassmann-lower-bound}
\leanverified{paper\_xi\_offline\_membership\_grassmann\_lower\_bound}
设 \(p\) 为素数，\(\FF_p\) 为有限域，取 \(n\in\NN\) 与 \(0\le r<n\)，并令 \(k:=n-r\)。
考虑如下离线成员判定任务：存在未知 \(k\)-维子空间 \(K\le \FF_p^n\)，发布比特串 \(\sigma\in\{0,1\}^\ell\) 作为验证参数，使得对任意查询向量 \(x\in\FF_p^n\)，验证者仅凭 \((\sigma,x)\) 即可零误差判定命题 \(x\in K\)。
则任何满足零误差成员判定的发布机制都必须满足
\[
\ell\ \ge\ \log_2\#\mathrm{Gr}(k,n;\FF_p)\ \ge\ k(n-k)\log_2 p\ =\ r(n-r)\log_2 p,
\]
其中 \(\mathrm{Gr}(k,n;\FF_p)\) 为 \(k\)-维子空间的 Grassmann 集。
\end{theorem}
\begin{proof}
若存在两不同子空间 \(K_1\neq K_2\) 对应同一 \(\sigma\)，取 \(x\in K_1\setminus K_2\) 则成员判定在输入 \((\sigma,x)\) 上必同时给出“真”与“假”，矛盾。
故发布映射必须在 Grassmann 集上单射，从而 \(\ell\ge \log_2\#\mathrm{Gr}(k,n;\FF_p)\)。
\par
由高斯二项式公式，
\[
\#\mathrm{Gr}(k,n;\FF_p)=\prod_{i=0}^{k-1}\frac{p^n-p^i}{p^k-p^i}
=\prod_{i=0}^{k-1}p^{n-k}\cdot\frac{1-p^{i-n}}{1-p^{i-k}}
\ge \prod_{i=0}^{k-1}p^{n-k}
=p^{k(n-k)}.
\]
取对数得第二个不等式。证毕。
\end{proof}

\begin{corollary}[子空间规范编码实现的近最短性：主阶达到 Grassmann 下界]\label{cor:xi-offline-membership-near-optimal-encoding}
在定理 \ref{thm:xi-offline-membership-grassmann-lower-bound} 的设定下，若将验证参数取为 \(K\) 的一个规范表示（例如给出 \(K\) 的行最简形基矩阵并附带主元位置编码），则存在零误差成员判定实现，其编码长度满足
\[
\ell\ \le\ r(n-r)\log_2 p\ +\ O(n\log n).
\]
因此在 \(p\) 固定而 \(n\to\infty\) 的主阶意义下，该编码达到 Grassmann 计数下界；子空间 \(K\) 在离线审计语义下构成近最短公钥。
\end{corollary}

% Shift-invariant saturated lattices and syndrome polynomial compression.

\paragraph{截断移位不变综合征子格的主理想结构}
在 Hankel 综合征的整数语义下，核约束子格自然携带截断移位作用，并由此形成一个 \(\ZZ[\lambda]\)-模结构。
在饱和（无挠）条件下，该模结构退化为截断多项式环中的主理想，从而把综合征子格压缩为一元首一本原多项式，并导出格运算、生成系与诱导移位谱指纹的完全刻画。

\begin{theorem}[饱和截断移位不变子格的主理想分类]\label{thm:xi-saturated-shift-invariant-lattice-principal-ideal}
取 \(n\ge 1\)。
将 \(\ZZ^n\) 与截断多项式模
\(
\ZZ[\lambda]_{<n}
\)
以系数向量同一，并令截断移位算子 \(\sigma:\ZZ^n\to\ZZ^n\) 对应于
\[
f(\lambda)\ \longmapsto\ \lambda f(\lambda)\bmod \lambda^n
\qquad\text{在}\qquad
R_n:=\ZZ[\lambda]/(\lambda^n)
\]
上的乘法作用。
设子格 \(L\subset\ZZ^n\) 满足：
\begin{enumerate}
  \item \(\sigma(L)\subseteq L\)；
  \item \(L\) 饱和，即 \(\ZZ^n/L\) 无挠。
\end{enumerate}
记 \(\Delta:=\rank_{\ZZ}L\) 且 \(m:=n-\Delta\)。
则存在唯一的首一本原多项式 \(P_L(\lambda)\in\ZZ[\lambda]\) 满足 \(\deg P_L=m\)，使得
\[
L=M_n(P_L):=\Bigl\{\operatorname{coeff}_{<n}\bigl(Q(\lambda)P_L(\lambda)\bigr):\ Q\in\ZZ[\lambda],\ \deg Q<n-\deg P_L\Bigr\}.
\]
其中 \(\operatorname{coeff}_{<n}\) 表示取模 \(\lambda^n\) 后的系数向量（等价于在 \(R_n\) 中的像）。
\end{theorem}
\begin{proof}
在识别 \(\ZZ^n\simeq R_n\) 下，条件 \(\sigma(L)\subseteq L\) 等价于 \(L\) 对 \(\lambda\) 的乘法封闭。
因 \(L\) 亦为 \(\ZZ\)-子模，遂 \(L\) 为环 \(R_n\) 的一个理想，记为 \(I\triangleleft R_n\)。
\par
由饱和性知 \(R_n/I\simeq \ZZ^{m}\) 为秩 \(m\) 的自由阿贝尔群，并由 \(\bar 1\) 在 \(\lambda\) 作用下生成。
令 \(m\) 为满足 \(\{\bar 1,\lambda\bar 1,\dots,\lambda^{m}\bar 1\}\) 在 \(\ZZ\) 上线性相关的最小整数，则
\(\{\bar 1,\dots,\lambda^{m-1}\bar 1\}\) 线性无关并构成 \(R_n/I\) 的一组 \(\ZZ\)-基。
因此存在唯一整数 \(p_0,\dots,p_{m-1}\) 使
\[
\lambda^{m}\bar 1=-\sum_{j=0}^{m-1}p_j\,\lambda^{j}\bar 1.
\]
定义首一多项式
\(
P(\lambda):=\lambda^{m}+\sum_{j=0}^{m-1}p_j\lambda^{j}
 \in\ZZ[\lambda]
\)。
则 \(P(\lambda)\bar 1=0\)，从而 \(P\in I\subset R_n\)。
\par
进一步，关系式给出 \(p_0\bar 1\in \lambda(R_n/I)\)，迭代得到 \(p_0^k\bar 1\in\lambda^{k}(R_n/I)\)。
取 \(k=n\) 并用 \(\lambda^n=0\) 得 \(p_0^n\bar 1=0\)。
由于 \(R_n/I\) 无挠且 \(\bar 1\neq 0\)，必有 \(p_0=0\)，从而 \(P(0)=0\) 与幂零截断相容。
\par
现证 \(I=(P)\) 于 \(R_n\)。
取任意 \(f\in I\)，选取次数 \(<n\) 的代表 \(\tilde f\in\ZZ[\lambda]\)。
对首一 \(P\) 作带余除法：
\(
\tilde f=QP+R
\)
其中 \(\deg R<m\)。
在 \(R_n\) 中有 \(f\equiv QP+R\)，故由 \(f\in I\) 且 \(P\in I\) 得 \(R\in I\)。
然而 \(R\bar 1=0\) 等价于 \(\sum_{j=0}^{m-1}r_j\lambda^j\bar 1=0\)，由基的线性无关性得 \(R=0\)。
于是 \(f\in(P)\)，从而 \(I\subseteq(P)\)；反向包含由 \(P\in I\) 立即成立，故 \(I=(P)\)。
\par
将理想 \((P)\subset R_n\) 写回系数向量即得 \(L=M_n(P)\)。
定义 \(P_L:=P\) 并将其归一化为首一本原元，得到存在性。
唯一性由主理想生成元的首一归一化唯一性给出。证毕。
\end{proof}
\leanverified{paper\_xi\_saturated\_shift\_invariant\_lattice\_principal\_ideal}

\begin{corollary}[饱和截断移位不变子格的格结构]\label{cor:xi-saturated-shift-invariant-lattice-gcd-lcm}
令 \(\mathcal L_n\) 为 \(\ZZ^n\) 中全体饱和且 \(\sigma\)-不变子格所成集合。
由定理 \ref{thm:xi-saturated-shift-invariant-lattice-principal-ideal}，映射
\[
\Psi:\mathcal L_n\to \mathcal P_n,\qquad L\mapsto P_L
\]
给出包含序与整除序之间的反序双射，其中 \(\mathcal P_n\) 为次数 \(<n\) 的首一本原整系数多项式集合。
更具体地，若 \(L_i=M_n(P_i)\)（\(i=1,2\)），则
\[
L_1\cap L_2\ =\ M_n(\mathrm{lcm}(P_1,P_2)),
\]
并且记 \(\operatorname{sat}(A)\) 为子群 \(A\subseteq\ZZ^n\) 的饱和闭包，则
\[
\operatorname{sat}(L_1+L_2)\ =\ M_n(\gcd(P_1,P_2)),
\]
其中 \(\gcd\) 取于 \(\QQ[\lambda]\) 的首一最大公因子并由 Gauss 引理提升为 \(\ZZ[\lambda]\) 的本原多项式。
\end{corollary}
\leanverified{paper\_xi\_saturated\_shift\_invariant\_lattice\_gcd\_lcm}
\begin{proof}
由 \(M_n(P)\) 与理想 \((P)\subset R_n\) 的对应，交 \(L_1\cap L_2\) 对应理想交 \((P_1)\cap(P_2)=(\mathrm{lcm}(P_1,P_2))\)。
\par
对和 \(L_1+L_2\) 亦为 \(\sigma\)-不变子群，其饱和闭包仍为 \(\sigma\)-不变。
设 \(\operatorname{sat}(L_1+L_2)=M_n(P)\)。
由包含反序对应得 \(P\mid P_1\) 且 \(P\mid P_2\)，故 \(P\mid \gcd(P_1,P_2)\)，从而 \(M_n(\gcd(P_1,P_2))\subseteq M_n(P)\)。
反向包含由 \(\gcd(P_1,P_2)\mid P_i\Rightarrow L_i\subseteq M_n(\gcd(P_1,P_2))\) 推出 \(L_1+L_2\subseteq M_n(\gcd(P_1,P_2))\)，再由后者饱和性得 \(\operatorname{sat}(L_1+L_2)\subseteq M_n(\gcd(P_1,P_2))\)。
合并即得。证毕。
\end{proof}

\begin{proposition}[主生成元的内禀刻画与自然生成系]\label{prop:xi-saturated-shift-invariant-lattice-intrinsic-generator}
设 \(L\in\mathcal L_n\)，令 \(P_L\) 为定理 \ref{thm:xi-saturated-shift-invariant-lattice-principal-ideal} 给出的唯一首一本原多项式，并记 \(\Delta:=\rank_{\ZZ}L\)、\(d:=\deg P_L=n-\Delta\)。
则：
\begin{enumerate}
  \item \(P_L\) 是 \(L\) 中次数最小的非零多项式，并且在 \(L\) 中该“最小次数首一元”唯一；
  \item \(L\cong \ZZ^{\Delta}\)，且
        \[
        L=\Bigl\langle \operatorname{coeff}_{<n}(\lambda^j P_L):\ 0\le j\le \Delta-1\Bigr\rangle_{\ZZ};
        \]
  \item 商群 \(\ZZ^n/L\) 为秩 \(d\) 的自由阿贝尔群，并由 \(\bar 1\) 在 \(\sigma\) 的作用下生成。
\end{enumerate}
\end{proposition}
\leanverified{paper\_xi\_syndrome\_certificate\_incompressibility}
\begin{proof}
由 \(L=M_n(P_L)\) 得 \(P_L\in L\)。
若 \(Q\in L\) 且 \(\deg Q<\deg P_L\)，则 \(Q=\operatorname{coeff}_{<n}(AP_L)\) 对某个 \(A\in\ZZ[\lambda]\) 成立；除非 \(A=0\)，否则 \(\deg(AP_L)\ge \deg P_L\)，故必有 \(Q=0\)。
这给出最小次数性质与唯一性。
第 (2)(3) 由 \(M_n(P_L)\) 的定义与定理 \ref{thm:xi-saturated-shift-invariant-lattice-principal-ideal} 的循环生成构造直接推出。证毕。
\leanverified{paper\_xi\_saturated\_shift\_invariant\_lattice\_intrinsic\_generator}
\end{proof}

\begin{proposition}[综合征证书的不可压缩度：从 \texorpdfstring{$\Delta$}{Delta} 向量到一元首一多项式的最小参数化]\label{prop:xi-syndrome-certificate-incompressibility}
设 \(L\in\mathcal L_n\) 且 \(L=M_n(P)\)，其中 \(P\) 首一本原且 \(\deg P=d<n\)。
则：
\begin{enumerate}
  \item 仅发送多项式 \(P\)（等价地发送其 \(d\) 个非首项系数）即可在验证端重建子格 \(L\) 的全部线性约束；
  \item 在允许任意 unimodular 基变换（即仅以 \(\GL_n(\ZZ)\)-不变量区分子格）的语义下，任何能够唯一决定该类子格的整系数不变量族，其自由参数数目至少为 \(d\)。
\end{enumerate}
\end{proposition}
\begin{proof}
第 (1) 点由命题 \ref{prop:xi-saturated-shift-invariant-lattice-intrinsic-generator}(2) 立即给出：\(P\) 唯一确定自然生成元族，从而唯一确定 \(L\)。
第 (2) 点：由定理 \ref{thm:xi-saturated-shift-invariant-lattice-principal-ideal} 的唯一性，首一本原多项式 \(P\mapsto M_n(P)\) 为单射。
因此若存在少于 \(d\) 个自由整数参数即可区分所有 \(M_n(P)\)，则必有不同 \(P_1\neq P_2\) 在该不变量族下不可区分，从而对应子格亦不可区分，矛盾。
故参数下界至少为 \(d\)。证毕。
\end{proof}

\begin{theorem}[综合征商模的谱指纹：由首一多项式给出诱导移位的有理标准型]\label{thm:xi-syndrome-quotient-module-rational-canonical-fingerprint}
令 \(L\in\mathcal L_n\) 且 \(L=M_n(P)\)，其中 \(P\) 为定理 \ref{thm:xi-saturated-shift-invariant-lattice-principal-ideal} 的首一本原多项式，\(\deg P=d\)。
记商模
\(
M:=\ZZ^n/L
\)，并令 \(\bar\sigma\) 为 \(\sigma\) 在 \(M\) 上诱导的端同态。
则：
\begin{enumerate}
  \item \(M\) 为秩 \(d\) 的自由阿贝尔群，且 \((M,\bar\sigma)\) 为循环 \(\ZZ[\lambda]\)-模，由 \(\bar 1\) 生成并满足 \(P(\bar\sigma)\bar 1=0\)；
  \item 在 \(\QQ\)-线性化后，\(\bar\sigma_{\QQ}:=\bar\sigma\otimes 1\) 的有理标准型由 \(P\) 在 \(\QQ[\lambda]\) 的素幂分解完全决定：若
        \(
        P=\prod_{i=1}^s p_i(\lambda)^{e_i}
        \)
        其中 \(p_i\) 两两互异不可约，则
        \[
        M\otimes_{\ZZ}\QQ\ \cong\ \bigoplus_{i=1}^{s}\QQ[\lambda]/(p_i^{e_i}),
        \]
        且各直和分量上的 \(\lambda\)-作用分别为 \(p_i^{e_i}\) 的伴随块（Frobenius/rational companion blocks）；
  \item \(\bar\sigma_{\QQ}\) 在 \(\CC\) 上半单当且仅当 \(P\) 在 \(\QQ[\lambda]\) 中平方自由；并且其在 \(\CC\) 上具有单谱当且仅当 \(P\) 具有互异复根。
\end{enumerate}
\end{theorem}
\begin{proof}
(1) 由定理 \ref{thm:xi-saturated-shift-invariant-lattice-principal-ideal}，\(\rank_{\ZZ}L=n-d\)，故 \(\rank_{\ZZ}M=d\)。
由构造，\(M\) 由 \(\bar 1\) 在 \(\bar\sigma\) 作用下生成，并且 \(P\in L\) 等价于 \(P(\bar\sigma)\bar 1=0\)。
\par
(2) 张量到 \(\QQ\) 后，\(\QQ[\lambda]\) 为 Euclid 域，有限生成循环 \(\QQ[\lambda]\)-模由其湮灭理想的素幂分解给出标准分解；该分解与 \(\bar\sigma_{\QQ}\) 的有理标准型一致。
\par
(3) \(\bar\sigma_{\QQ}\) 的极小多项式为 \(P\) 的平方自由部分，故在代数闭域上半单当且仅当极小多项式平方自由；单谱等价于特征多项式（此处为 \(P\)）具有互异根。证毕。
\end{proof}

\endinput

\begin{proof}
行最简形给出子空间的规范代表。
其自由参数可组织为一个 \(k\times(n-k)\) 的矩阵（自由列上的系数），共 \(k(n-k)=r(n-r)\) 个 \(\FF_p\) 元素；其存储成本为 \(r(n-r)\log_2 p\) 比特。
主元列集合的编码可在 \(O(n\log n)\) 比特内完成。
成员判定可由该规范代表构造线性方程组并作确定性消元实现，且零误差。证毕。
\end{proof}

\endinput
